\section{Reglas empíricas de los Compound AI Systems}

\subsection{Comparación de los patrones principales}

\begin{table}[H]
\centering
\caption{Patrones de construcción de los Compound AI Systems}
\begin{tabular}{p{0.2\textwidth}p{0.26\textwidth}p{0.22\textwidth}p{0.24\textwidth}}
\toprule
Patrón & Composición & Ventaja & Reto de ingeniería \\
\midrule
RAG & Recuperación + 1 llamada a LLM & Implementación sencilla & Carece de capacidad de razonamiento encadenado \\
ReAct & Ciclo de razonamiento → herramienta → observación & Adecuado para tareas de varios pasos & Requiere una lógica de control minuciosa \\
DSPy & Estructura declarativa + optimizador & Programable y transferible & Espacio de búsqueda grande, el ajuste consume mucho tiempo \\
Autonomous Agent & Colaboración de varios Agent + MCP & Flujos de trabajo paralelos automatizados & Requiere monitoreo fuerte y barreras de seguridad \\
\bottomrule
\end{tabular}
\end{table}

\begin{knowledgebox}{Reglas empíricas}
El diseño de los sistemas de IA compuestos sigue tres principios: 1) los componentes necesitan una firma clara para poder componerse; 2) la estructura debe verificarse con metric / log; 3) el proceso de optimización debe ser reproducible, o de lo contrario el sistema es propenso al drift.
\end{knowledgebox}

\subsection{Firma y composición}

Signature define la forma de entrada/salida de un módulo; DSPy usa la verificación de signature para evitar combinaciones no válidas. Cada vez que se amplía el pipeline, se debe insertar un verificador de tipos y generar automáticamente pruebas de interacción, para evitar que los módulos posteriores reciban valores anómalos.

\subsection{Reutilización de patrones}

Los patrones comunes de los sistemas compuestos (como retriever → reasoner, planner → executor) pueden encapsularse como reusable module. Una vez que la firma es consistente, pueden reutilizarse de manera recurrente en distintos pipeline, lo que reduce considerablemente el costo de copiar prompts a mano.

\begin{knowledgebox}{La reutilización exige firmas consistentes}
Ni el optimizador más potente puede proteger las combinaciones que no coinciden. Si un módulo ``state → action'' se inserta por error en un pipeline ``question → answer'', se generarán continuamente invalid token en la Evidence Matrix. El verificador de firmas debe lanzar un error ya en la etapa de compile.
\end{knowledgebox}

\begin{warningbox}{No trates el pipeline como una caja negra}
En un pipeline que incluye recuperación, razonamiento, herramientas y validación, un error en cualquier paso puede propagarse, así que colocar guard rails (aserciones, verificación de tipos, fallback) en cada límite de módulo es más importante que optimizar el prompt.
\end{warningbox}

\subsection{Resumen de la sección}

La confiabilidad de los Compound AI Systems proviene de interfaces claras, la verificación de que los componentes puedan combinarse y un proceso de optimización observable, no solo de depender de un LLM más potente.
